\section{Arithmetic coordinates and common local relations}\label{sec:arithmetic}

We construct the algebraic coordinates and common local equations needed
for \Cref{thm:interpolation}.  The coordinates will have controlled global
height and much smaller norms on the center disks.  A weighted proximity
estimate then ensures that one common system of equations retains enough
smallness of the local parameter.  Throughout this section the curve, its
finite cover, its algebraic frames, and the center are fixed.  In
particular the conclusions apply after any one fixed auxiliary level
cover.  No condition on the reduction type of the center is imposed.

Write $\pi:\mathcal A\to S$ for the abelian scheme and put
$\mathcal H=R^1\pi_*\Omega^\bullet_{\mathcal A/S}$.  We use the
operator realization of $\Vcal\subset\End(\mathcal H)$ from
\Cref{lem:qm-plane}.  On actual endomorphisms, $\Tr$ denotes the trace
on first cohomology, so that $\Tr(1)=4$.  As in
\eqref{eq:quadratic-system}, $\pair{u}{v}=\tfrac14\Tr(uv)$.
For an endomorphism $u$, set
\[
 q_\lambda(u)=\Tr(uu^\dagger),
\]
where $\dagger$ is the principal Rosati involution.  This is a positive
definite integral quadratic form on $\End(A)$.

For a number field $F$ containing coordinates under consideration, our
weights are $w_v=[F_v:\Q_v]/[F:\Q]$, with the usual unsquared absolute
values.  Thus $\ha{z}=\sum_v w_v\log\max(1,\norm{z}_v)$ for an
affine tuple $z$, where the norm is the coordinate maximum.  The same sums
can be written over embeddings into $\C$ or $\overline{\Q}_p$, giving
each embedding weight $1/[F:\Q]$.  If a selected set of embeddings is
extended to a larger field, we always take \emph{all} extensions.  Their
weights sum to the old weight.  This convention also permits grouping
embeddings in a normal closure without charging its degree to any estimate.

After omitting finitely many points, fix a rational frame of $\Vcal$
regular at $t_0$, affine coordinates $z$ on a neighborhood of $t_0$, and
$x\in K(S)$ with a simple zero at $t_0$.  The inverse branch through
$t_0$ is denoted $z=z(X)$, with $z(0)=z(t_0)$.  Shrink the neighborhood
so that these data and the connection are regular there.
Choose two independent integral trace-free Rosati-symmetric endomorphisms
of $A_{t_0}$ as in \Cref{lem:qm-plane}, enlarge $K$ once to define them,
and denote their de Rham coordinate pair by $B=(B_1,B_2)$.
For ordered vector pairs $\mathbf v=(v_1,v_2)$ and
$\mathbf w=(w_1,w_2)$, write
$\pair{\mathbf v}{\mathbf w}=(\pair{v_i}{w_j})_{i,j=1}^2$.
Put
$a=x(t)\ne0$, $h=\max\{2,h_H(t)\}$ for a fixed nonnegative
ample logarithmic height on $\overline S$, and
$d=\max\{2,[K(t):K]\}$.
All constants below depend only on the fixed data.

\subsection{Integral endomorphisms and de Rham coordinates}

\begin{lemma}[Potential good reduction]\label{lem:potential-good}
An abelian surface whose geometric endomorphism algebra is a quaternion
division algebra over $\Q$ has potentially good reduction at every finite
place.
\end{lemma}

\begin{proof}
Fix a finite place.  The semistable reduction theorem
\cite[Section~7.4, Theorem~1]{BLR1990}, followed by a further finite
extension, gives semistable reduction and defines all geometric
endomorphisms.  They extend to the N\'eron model and preserve
the torus in its connected special fiber.  If that torus has positive
rank $r$, its rational character space is a nonzero unital module, on
one side or the other, for the division algebra
$D=\End^0(A_{\overline K})$.  Every finite dimensional module over a
division algebra is free over it, so its rational dimension is a
positive multiple of $\dim_\Q D=4$.  This contradicts $r\le2$.
The toric rank is therefore zero, and the semistable model is an
abelian scheme.  The extension used here is local; no uniform bound on
its degree is needed for an absolute height computation.
\end{proof}

\begin{proposition}[Small labeled integral vectors]\label{prop:small-vectors}
There are constants $c,C>0$ with the following property.  For every
target $t$, an extension $L/K(t)$ of degree at most $C$ defines all
geometric endomorphisms of $A_t$.  There are two independent
endomorphisms $u_1,u_2\in\End_L(A_t)$ which are symmetric and trace
free, span $E_t$, and satisfy
\begin{equation}\label{eq:rosati-size}
 q_{\lambda_t}(u_j)\le C(dh)^c\qquad(j=1,2).
\end{equation}
The corresponding labeled vectors are denoted $P=(P_1,P_2)$ in de Rham
coordinates.  Conjugating a labeled target always means conjugating
these endomorphisms as well.  This pair need not generate the saturated
integral cohomology lattice of $E_t$.
\end{proposition}

\begin{proof}
The Galois action on the geometric endomorphism lattice has finite
image in $\mathrm{GL}_r(\Z)$, with $r\le16$.  Indeed a finite set of
lattice generators is defined over a finite extension.  A finite
subgroup of $\mathrm{GL}_r(\Z)$ injects into
$\mathrm{GL}_r(\mathbf F_3)$, because the kernel of reduction modulo
$3$ is torsion free.  Its order is consequently bounded in terms of
$r$ alone.  The fixed field of the kernel gives the asserted $L$.

For a principally polarized abelian variety of fixed dimension, the
endomorphism estimate of Masser--W\"ustholz gives a rational spanning
set of integral endomorphisms whose Rosati lengths are bounded by a
fixed power of
\[
 \max\{2,[L:\Q],h_F(A_t)\}.
\]
Here polynomial dependence on the field degree is essential.  It is
explicit in the Proposition of \cite[Section~4]{MW1994}, with its
parameter $m$ equal to the dimension of the abelian variety; the
alternative involving a smaller image is then zero.  Its lattice
Lemma~2.1 also gives a basis with the same kind of bound.  See also
\cite[Theorem~6.6 and Section~6.7]{DawOrr2022} for the polynomial
dependence in this application.  The theta-height comparison
\cite[Theorem~1.1 and Corollary~1.3]{Pazuki2012}, on one fixed theta
level, and functoriality of heights on a fixed curve give
\[
 h_F(A_t)\le C h.
\]
The degree of passage to that fixed level is bounded.

For clarity, the projection to the required plane can be made without
target-dependent denominators.  The map
\[
 u\longmapsto 4(u+u^\dagger)-\Tr(u+u^\dagger)1
\]
preserves integrality, has symmetric trace-free image, and increases
the Rosati norm by at most a fixed factor.  By
\Cref{lem:qm-plane} its rational image has dimension two.  Apply this
map to the small spanning set and select two independent images.
This proves \eqref{eq:rosati-size}.  The norm and trace assertions are
unchanged under field conjugacy: they are traces of actual
endomorphisms and their Rosati adjoints.
\end{proof}

The next distinction is useful.  A rational frame can have large
denominators at a moving point, even when the endomorphism is integral.
Their global cost is linear in $h$.  On fixed sufficiently small
neighborhoods of the center the frames are uniformly controlled, and
the restricted cost is much smaller.

\begin{proposition}[Coordinate heights and local norms]\label{prop:dr-size}
Choose the vectors of \Cref{prop:small-vectors}.  Their affine
coordinates in the fixed frame satisfy
\begin{equation}\label{eq:dr-global}
 \ha{z(t),a,P}\le C(h+\log d).
\end{equation}
There are center neighborhoods at every place, with a single integral
description outside a fixed finite set of rational primes, such that
for any selected weighted set $\mathcal V$ on the actual center branches,
\begin{equation}\label{eq:dr-local}
 \sum_{v\in\mathcal V}w_v
       \log\max(1,\norm{z(t),P}_v)
       \le C(1+\log h+\log d).
\end{equation}
These statements hold for the finitely many conjugate fixed charts.
They also hold after a bounded denominator and a polarized isogeny
of any one fixed multiplier are applied to the vectors.
\end{proposition}

\begin{proof}
\emph{Global finite-place bounds.}
We first bound the denominators of integral endomorphisms in the rational
frame.  Potential good reduction supplies an integral de Rham lattice at
each target; a fixed ideal measures the denominators of the frame and
its dual relative to that lattice.  Pass,
for this argument only, to a fixed common cover carrying two coprime
fine levels $N_1,N_2\ge3$.  The integral fine moduli scheme
$\Acal_2(N_i)$ is quasi-projective over $\Ocal_K[1/N_i]$
\cite[Chapter~I, Theorem~1.4 and Section~1.7]{Chai1985}.
Embed it in projective space over this ring, take its scheme-theoretic
closure, and then take the closure of the graph of the curve map in
the product with the fixed projective curve model.
Denote this proper graph
model by $X_i$.  Take its projective closure over $\Ocal_K$ when
using global model heights.  The inverse image $U_i$ of the abelian interior
carries the integral vector bundle $\mathcal H_i$ of first de Rham
cohomology and its dual.  At a prime $p\nmid N_i$, a target acquires
good reduction by \Cref{lem:potential-good}; its prime-to-$p$ level
extends after a finite local extension.  Properness extends the
point to $X_i$.  The good-reduction abelian scheme with prime-to-$p$
level also supplies a valuation-ring point of the abelian interior.
Its map to the projective closure agrees with the graph map on the
generic fiber, hence everywhere by separatedness.  The moduli image
of the extended graph point therefore lies in the abelian interior.
Consequently the extended point belongs to $U_i$.

This last assertion concerns the moduli image.  The point is allowed
to reduce above a boundary point of the original curve.  Thus we use
$U_i$, rather than a preassigned integral open of the curve, to
compare lattices.  The two choices of $i$ cover every residue
characteristic.

Choose a rational frame $e_1,\ldots,e_4$ of $\mathcal H$ such that it
and its dual $e^1,\ldots,e^4$ are regular at $t_0$.  Form on $U_i$ the
coherent denominator ideals
\[
 \mathcal I_+=\{f:f e_j\in\mathcal H_i\text{ for all }j\},
 \qquad
 \mathcal I_-=\{f:f e^j\in\mathcal H_i^\vee\text{ for all }j\}.
\]
These are coherent colon ideals, and extend to coherent ideals on
the noetherian projective model $X_i$
\cite[Lemma~28.23.1(2), Tag~01PE]{StacksProject}.  Their product, further
multiplied by a denominator ideal for the fixed rational conversion
from operator entries to our frame of $\Vcal$, gives a fixed ideal
$\mathcal J_i$.  Multiplication by $\mathcal J_i$ makes all the
specified sections integral on $U_i$.  Remove its finitely many
generic zeros from the curve.  No duality assertion about arbitrary
coherent extensions across the boundary is required.
For a lifted target over a valuation ring $R$, write
$\mathcal J_i(t)R=(\eta)$ and
$\Lambda=H^1_{\mathrm{dR}}(\mathcal A_t/R)$.  The frame vectors
belong to $\eta^{-1}\Lambda$ and their duals belong to
$\eta^{-1}\Lambda^\vee$.  An actual integral endomorphism preserves
$\Lambda$: it extends to the abelian scheme
\cite[Section~1.2, Proposition~8]{BLR1990}, and de Rham cohomology
is functorial.  Each of its matrix entries is therefore in
$\eta^{-2}R$.  The extra fixed conversion of coordinates gives,
for some fixed integer $b$,
\begin{equation}\label{eq:pole-ideal}
 \log\max(1,\norm{P}_v)
       \le b\lambda_{\mathcal J_i,v}(t)+c_v
       \quad(v\nmid N_i),
\end{equation}
where $c_v=0$ outside finitely many fixed primes.

Here is explicitly why the right side has the claimed global cost.
For a sufficiently ample fixed line bundle $\mathcal L$ on $X_i$,
$\mathcal J_i\mathcal L^n$ is generated by finitely many fixed
sections $s_0,\ldots,s_r$
\cite[Proposition~28.27.13(7), Tag~01Q3]{StacksProject}.
Choose $s_k$ with $s_k(t)\ne0$.
With the model norm at a finite place,
$\lambda_{\mathcal J_i,v}(t)\le-\log\norm{s_k(t)}_v$.
The product formula for this nonzero element of the metrized line
$\mathcal L_t^n$ gives
\[
 \sum_{v\nmid\infty}w_v\lambda_{\mathcal J_i,v}(t)
 \le n h_{\mathcal L}(t)
       +\sum_{v\mid\infty}w_v\log\norm{s_k(t)}_v
 \le n h_{\mathcal L}(t)+C.
\]
The last bound uses the bounded sup norms of the finitely many
fixed sections on the fixed projective complex model.  The finite
ideal heights are nonnegative, so this also bounds any restricted
finite sum, such as the sum over $p\nmid N_i$.
Every generic projective curve and map used here is fixed, so
$h_{\mathcal L}(t)=O(h)$.  Equation \eqref{eq:pole-ideal}, summed
using the two models, bounds all finite denominator contributions
by $O(h)$, independently of the local extension used for good
reduction.

\emph{Global archimedean bounds.}
At an archimedean place the Hodge norm of an actual endomorphism is
$q_{\lambda_t}^{1/2}$, up to a fixed convention.  In a rational de
Rham frame, the norms of the frame and dual have at most a fixed
power of the inverse distance to the finitely many omitted points,
times powers of its logarithm.  The algebraic Gauss--Manin connection
is regular singular by \cite[Theorem~14.1]{Katz1970}; this property
passes to its duals, tensor products and subquotients
\cite[Proposition~11.3 and Remark~11.3.4]{Katz1970}.
After a fixed ramified cover at each puncture, compare its algebraic
logarithmic lattice with the canonical extension of the variation.
Both frames and their inverses differ by meromorphic matrices with
finite pole orders.  Schmid's norm estimates for the variation and
its dual \cite[Theorems~4.9 and~6.6$'$]{Schmid1973} then give the
asserted power and logarithmic bounds for the rational frame and its
dual.  The sum of the positive logarithms of these factors is
$O(h)$ by the local height inequalities for the fixed omitted
divisor.  The contribution of \eqref{eq:rosati-size} is
$O(\log h+\log d)$.  Affine base coordinates and $a$ are fixed
rational functions and have height $O(h)$.  This proves
\eqref{eq:dr-global}.

\emph{Bounds on the center disks.}
For \eqref{eq:dr-local}, outside finitely many primes the family,
parameter chart, frame, and dual are integral near the center,
with the frame invertible.  All integral endomorphism coordinates
then have norm at most one on that residue neighborhood.  At a fixed exceptional prime $p$, extend the local coefficient
field once so that the fixed center has good reduction and a
chosen full $M$-level structure, where $M\ge3$ and $p\nmid M$.
Adding this level to the original characteristic-zero curve
gives a finite \'etale cover; at the chosen lift of the center,
the $p$-adic inverse function theorem gives an analytic section
over a sufficiently small center disk.  Every sufficiently
close target on the original branch therefore has a lift on
this section after coefficient extension, including when the
original level is divisible by $p$.  Forgetting the original
level maps the lifted disk to the prime-to-$p$ fine moduli
scheme.  On a smaller closed disk its coordinates lie in an
integral open about the fixed center and reduce to the same
special point.  An integral de Rham frame on this moduli open
pulls back to a frame whose change-of-basis matrix with the
original frame, and its inverse, are analytic and bounded on
the disk.  Thus both the finite-place comparison cost and the
base-coordinate norms are bounded by fixed local constants.

At infinity choose compact
disks within the regular charts; the frame and dual Hodge norms
are bounded there.  Only $O(\log h+\log d)$ remains from the
Rosati norm.  The sum of weights above a fixed rational prime is
one, as is the sum of archimedean weights, which proves the
restricted sum.

Finally a polarized isogeny of fixed multiplier carries Rosati
norms isometrically under conjugation on operators.  Its inverse
introduces only a bounded denominator into an actual
endomorphism.  This changes the finite estimates at finitely many
fixed primes by fixed constants, and proves the last assertion.
\end{proof}

\subsection{Horizontal comparison on good-reduction disks}

We next express pairings between centered and target endomorphisms in
one cohomology space.  At finite places the common special fiber supplies
that space, and functorial comparison must apply to each individual
endomorphism even when it does not extend over the base curve.
Let $Y(X)$ be the rank-five horizontal transport \emph{back} to the
center, normalized by $Y(0)=I$.  It is obtained from the corresponding
transport on $\mathcal H$ by the operator construction.  Its entries
are ordinary-point Taylor series over $K$.

\begin{proposition}[Functorial good-reduction comparison]\label{prop:good-comparison}
At a finite place, suppose a point $t$ is on a sufficiently small
actual center branch.  After coefficient extension the two
polarized abelian schemes have a common special fiber $\overline A$.
More precisely, over a finite local extension $F$ with residue
field $k$, put $F_0=W(k)[1/p]$ and
$D=H^1_{\mathrm{cris}}(\overline A/W(k))[1/p]$.  There are rational
comparison isomorphisms
\[
 c_u:H^1_{\mathrm{dR}}(A_u/F)\longrightarrow D\otimes_{F_0}F,
 \qquad u=t_0,t,
\]
whose transition $c_{t_0}^{-1}c_t$ is horizontal transport.
They are compatible with morphisms of either lifted fiber and
their specializations, including morphisms which do not extend
over the base curve.  Ramification of the coefficient field is
allowed.

In particular, if $\beta_i$ and $\alpha_j$ are the specializations
of the centered and target endomorphisms, then
\begin{equation}\label{eq:cross-trace}
 \pair{B_i}{Y(a)P_j}=\tfrac14\Tr(\beta_i\alpha_j)\in\Q
       \qquad(1\le i,j\le2).
\end{equation}
The cross pairings belong to $\tfrac14\Z$, and their ordinary
absolute values are at most $C(dh)^c$.  At archimedean places the
same assertions hold with integral Betti transport in place of
specialization.
\end{proposition}

\begin{proof}
Work on a smooth integral prime-to-$p$ fine moduli chart
$\mathfrak M$ over a Witt base, and let $z$ be the common
special point.  The relative first crystalline cohomology of
the universal abelian scheme defines a convergent isocrystal
$\mathcal E$ on the special fiber of this chart.  Its value
at $z$ is the crystalline cohomology of the actual abelian
variety $\overline A$ over $z$.  The canonical realization of
$\mathcal E$ on any smooth proper lifting is its relative
de Rham cohomology with Gauss--Manin connection; see
\cite[Theorems~4.16 and~4.26]{MaulikPoonen2012}, the latter
using \cite[Theorem~3.8.2]{Ogus1984}.

Here is the common evaluation explicitly.  Over the local
field $F$ in the statement, put
$\mathfrak T=\operatorname{Spf}\Ocal_F\langle Z\rangle$
and substitute $X=aZ$ into the chosen inverse branch.  Shrinking the
fixed center disks ensures that the resulting chart coordinates belong
to $\Ocal_F\langle Z\rangle$ and differ from their centered values by
elements of $\mathfrak m_F\Ocal_F\langle Z\rangle$, where
$\mathfrak m_F$ is the maximal ideal of $\Ocal_F$.  Thus the map to
$\mathfrak M$ reduces constantly to $z$ on
$(\mathfrak T\bmod p)_{\mathrm{red}}=\operatorname{Spec}k[Z]$.
This remains true when $F$ is ramified and $a\notin p\Ocal_F$:
the image of $\mathfrak m_F$ in $\Ocal_F/p$ is nilpotent.
Consequently $\mathfrak T$ is an enlargement of $z$, and
the crystal transition for its structural map to the point
enlargement gives a canonical horizontal identification
\[
 \mathcal E_{\mathfrak T}
   \simeq D\otimes_{F_0}\Ocal_{\mathfrak T}[1/p].
\]
This is the constant-reduction construction of
\cite[Definitions~4.10--4.13 and Section~4.7]{MaulikPoonen2012}.

At $Z=0$ and $Z=1$, the relative comparison gives precisely
the two maps $c_{t_0}$ and $c_t$ of the statement.  Their
compatibility with restriction to a fiber identifies them
with the functorial absolute crystalline--de Rham
comparisons; an explicit absolute statement with arbitrary
ramification is
\cite[Appendix~B, Proposition~B.4]{KubrakPrikhodko2022}.
Proper formal GAGA identifies completed and ordinary
de Rham cohomology.  An endomorphism of either individual
fiber extends over its good-reduction abelian scheme and
is carried to its actual specialization on $\overline A$.
This argument needs only that individual endomorphism,
without an extension to neighboring fibers.  The common
fiber is $\overline A$ itself, with the same specialization
maps and prime-to-$p$ markings as in smooth proper base
change.  The construction does not introduce an auxiliary
isogeny between special fibers.

The crystal identification is horizontal and is the
identity when the two evaluations coincide.  Therefore
$c_{t_0}^{-1}c_t$ is the identity-normalized horizontal
transport, and uniqueness for the ordinary differential
equation identifies its operator realization with $Y(a)$.
The resulting relative matrix identity is also discussed
in \cite[Theorem~3.4 and Remark~3.6]{PapasI2025}.

Outside finitely many
fixed primes our curve chart and connection are integral; the
Taylor comparison converges, including over ramified extensions,
whenever
\begin{equation}\label{eq:factorial-disk}
 -\log|a|_v>\frac{2\log p}{p-1}.
\end{equation}
Indeed the divided-power terms have denominators dividing the
corresponding factorial.  At the finitely many other primes,
\Cref{lem:potential-good} for the fixed center and a
prime-to-characteristic moduli chart give the same assertion on
a fixed smaller disk.

Functoriality identifies both operator realizations with the
operators $\beta_i,\alpha_j$ on the same crystalline vector
space, proving \eqref{eq:cross-trace}.  The characteristic
polynomial of an endomorphism of an abelian variety has integral
coefficients and is independent of the Weil cohomology: it is
characterized by the degree polynomial $\deg(n-u)$
\cite[Section~12]{Milne1986}.  For the crystalline realization used
here, the cohomological assertion can be seen directly.  Rational
crystalline cohomology is a Weil cohomology
\cite[Remark~4.15]{MaulikPoonen2012}.  Cup product identifies the
cohomology of the abelian special fiber with the exterior algebra on
$H^1$: comparison with its good lift reduces this to the corresponding
statement for a complex abelian variety.  Thus $[n]-u$ acts on top
cohomology by $\det(n-u\mid H^1)$; the degree formula identifies that
action with $\deg([n]-u)$.  Equality for all sufficiently large integers
$n$ identifies the characteristic polynomials, and hence their traces.
This applies
also to $\beta_i\alpha_j$, so its trace is integral and the pairing
in \eqref{eq:cross-trace} belongs to $\tfrac14\Z$.  Specialization
preserves the polarization, Rosati involution, and these trace
polynomials.  Positivity of the Rosati form in characteristic
$p$ \cite[Theorem~17.3]{Milne1986} gives
\[
 |\Tr(\beta_i\alpha_j)|^2
 \le \Tr(\beta_i^2)\Tr(\alpha_j^2)
 \le C(dh)^c.
\]
At infinity parallel transport identifies integral first
cohomology lattices, so the cross pairing again belongs to
$\tfrac14\Z$.  On
the compact center disk the Hodge norms at the two points are
uniformly comparable under this transport.  Cauchy--Schwarz for
the resulting positive norm and \eqref{eq:rosati-size} give the
same size bound.  Notice that the two transported operators need
not both be Hodge endomorphisms at the center in this last argument.
\end{proof}

\subsection{Smallness on the actual center branches}

\begin{lemma}[Proximity and the factorial cutoff]\label{lem:proximity}
Use the fixed center disks of \Cref{prop:dr-size,prop:good-comparison},
and at every remaining prime impose \eqref{eq:factorial-disk}.
Let $\mathcal V_0(t)$ be all weighted embeddings at which $t$
belongs to these actual branches.  Then
\begin{equation}\label{eq:proximity-sum}
 \sum_{v\in\mathcal V_0(t)}w_v\log(1/|a|_v)
       \ge c h-C\bigl(1+\log^2(dh)\bigr).
\end{equation}
In particular the sum is at least $c h/2$ if $h\ge Q_0(d)$
for a fixed polynomial $Q_0$.
\end{lemma}

\begin{proof}
Use the effective divisor $D_0=[t_0]$ on the smooth completion
$\overline S$.  Since it has positive degree, some fixed multiple
of $D_0$ minus the ample divisor defining $h$ is ample.  The
height of an ample divisor is bounded below, and hence
\[
 h_{D_0}(t)\ge c h-C.
\]
Take local Weil functions for this divisor.  On the actual
branch through $t_0$ they equal $-\log|x(t)|_v$ up to fixed
bounded errors.  Outside the chosen branch neighborhoods their
positive contributions are bounded; the error is zero at all
but finitely many primes.  At an unexceptional prime the
integral parameter chart is \'etale at the center, so the branch
is precisely the Henselian inverse branch with that reduction.
Other zeros of the rational function $x$ are not counted.
Before the factorial cutoff, the weighted proximity sum is
therefore at least $ch-C$.

Here is a uniform bound for the loss caused by that cutoff.  Put
$D=[K(a):\Q]\le Cd$ and $H=\ha{a}\le Ch$.  If a rational prime
$p$ has any place where $|a|_v<1$, it contributes at least
$(\log p)/D$ to $\ha{a^{-1}}=H$: writing a local valuation in
integral units cancels its ramification index against the
absolute weight.  For the set $\mathcal P(a)$ of these primes,
\begin{equation}\label{eq:prime-support}
 \sum_{p\in\mathcal P(a)}\log p\le DH,
 \qquad \#\mathcal P(a)\le DH/\log2.
\end{equation}
Write $\alpha_p=\log p/(p-1)$.  At places over $p$ discarded
by \eqref{eq:factorial-disk}, the total weighted loss is at
most $2\alpha_p$, since their weights sum to at most one.
For any set of $N$ primes,
\[
 \sum\alpha_p
 \le C\sum_{k=1}^{N}\frac{\log(k+1)}k
 \le C\bigl(1+\log^2(N+1)\bigr).
\]
Here $\log x/(x-1)$ decreases for $x>1$, and the $k$th
smallest prime is at least $k+1$.  Apply this with
\eqref{eq:prime-support}.  Restricting the finitely many
exceptional and archimedean disks costs only a fixed further
constant.  This proves \eqref{eq:proximity-sum}; its last
assertion follows by elementary absorption of the logarithmic
term for $h\ge Q_0(d)$.
\end{proof}

\subsection{Grouping the equations}

Include in the fixed exceptional set all denominator primes of the
centered vectors and the trace pairing, in all fixed conjugate charts.
At infinity and at the finitely many exceptional primes, only $O(d)$
embeddings of the labeled field occur, so we can keep the rational
cross-trace matrix as a coefficient of each group's equations.  The
number of remaining primes can grow with the target.  There we compare
a trace equation with its image under a lift of residue Frobenius.
The rational trace is fixed, so subtraction eliminates it; only $O(d)$
conjugates of the complete labeled tuple can occur.  Retaining the labels
also records Frobenius's action on the actual endomorphisms, as needed
in \Cref{sec:height}.

\begin{proposition}[Common relations with controlled coefficients]
\label{prop:common-relations}
Suppose $h\ge Q_0(d)$, increasing $Q_0$ if necessary.  The
proximity embeddings can be partitioned into at most $Q_1(d)$
groups, and one group $\mathcal V$ has
\begin{equation}\label{eq:group-smallness}
 S_{\mathrm{small}}=
 \sum_{v\in\mathcal V}w_v\log(1/|a|_v)
       \ge h/Q_2(d).
\end{equation}
On each group there is one fixed list of analytic equations,
evaluated at one algebraic tuple $(a,U)$, of one of the following
forms:
\begin{align}
 \pair{B}{Y(a)P}&=T,
       &&\text{with }T\in\mathrm{M}_2(\Q),
       \label{eq:constant-relations}\\
 \pair{B}{Y(a)P}&=\pair{B'}{Y'(aR)P'},
       &&R=b/a,\quad b=x'(t'),
       \label{eq:paired-relations}
\end{align}
together with the actual branch equations
\begin{equation}\label{eq:chart-relations}
 z-z(a)=0,
 \qquad z'-z'(aR)=0
 \quad\text{when a second chart is present}.
\end{equation}
The primed data are one simultaneous field conjugate of the
complete labeled unprimed data.  In \eqref{eq:paired-relations}
all places are finite and $|R|_v=1$ on the group.  The Frobenius
of the common reduction transports each of the two labeled
pairs to its primed pair.

The tuple has degree at most $Q_3(d)$ over $K$, and
\begin{equation}\label{eq:common-size}
 \ha{a,U}\le Q_3(d)h,
 \qquad
 \sum_{v\in\mathcal V}w_v\log\max(1,\norm{U}_v)
       \le Q_3(d)(1+\log h).
\end{equation}
Writing the common list as
$E(a,U)=\sum_{j\ge0}a^j e_j(U)$, it has fixed length and
\begin{equation}\label{eq:coefficient-size}
 \begin{aligned}
 \deg e_j&\le C(1+j) &&(j\ge0),\\
 \ha{\text{affine coefficient tuple of }(e_0,\ldots,e_l)}
       &\le Q_4(l,d)(1+\log h) &&(l\ge0).
 \end{aligned}
\end{equation}
There are $b_v\ge0$ with
\begin{align}
 \left\|\sum_{j\ge l}a^j e_j(U)\right\|_v
       &\le |a|_v^l\exp\bigl((1+l)b_v\bigr)
             &&(l\ge0),\label{eq:arithmetic-tail}\\
 \sum_{v\in\mathcal V}w_v b_v
       &\le Q_5(d)(1+\log^2 h).
             \label{eq:tail-sum}
\end{align}
The series are formed from fixed ordinary-point algebraic
connection systems in $a$ and $aR$, and fixed algebraic chart
series, by polynomial operations of fixed degree.
\end{proposition}

\begin{proof}
First choose a field $L$ defining the complete labeled tuple:
the point, its two endomorphisms, and the fixed curve, center,
and frames.  By \Cref{prop:small-vectors}, $[L:\Q]\le Cd$.
Let $M$ be a normal closure.  We do not bound $[M:\Q]$.
Compute weights by extending every selected embedding to
\emph{all} embeddings of $M$; every old weighted sum is preserved.

At infinity or a fixed exceptional prime, the four entries of
the cross trace are the rational constants of
\Cref{prop:good-comparison}.  At any one such place type, an
embedding of the labeled field $L$ into its local algebraic
closure determines the Taylor evaluation and the constant.
There are at most $[L:\Q]$ possibilities.  Group by these
embeddings, by the finitely many fixed charts, and then by the
constant matrix.  This gives $O(d)$ groups of the form
\eqref{eq:constant-relations}, with
$\ha{T}\le C(1+\log d+\log h)$.
We count the constants actually obtained from these embeddings;
we do not count all rational matrices of that height.

At an unexceptional prime $p$, for each embedding
$\rho:M\hookrightarrow\overline\Q_p$, choose
$\sigma\in\operatorname{Gal}(\overline\Q_p/\Q_p)$ inducing the $p$th-power automorphism
of $\overline{\mathbf F}_p$.  Apply $\sigma$ to
\eqref{eq:cross-trace}.  The trace on the right is rational,
and is fixed by $\sigma$.  Thus
\[
 \pair{\rho B}{Y^\rho(\rho a)\rho P}
  =\pair{\sigma\rho B}
     {Y^{\sigma\rho}(\sigma\rho a)\sigma\rho P}.
\]
On $L$, the embedding $\sigma\rho$ equals $\rho\tau$ for
one of the at most $[L:\Q]$ embeddings $\tau:L\to M$.
Partition by $\tau$ and the fixed conjugate charts.  Define
$(B',P',t',b)=\tau(B,P,t,a)$ in this group.  This gives the
same analytic equation \eqref{eq:paired-relations} at every
embedding of the group.  The compositum $L\tau(L)$ has
degree at most $[L:\Q]^2$, so the field defining the
coordinates of this equation has polynomial degree even
though $M$ need not.

A local automorphism preserves the absolute value, whence
$|b|_v=|a|_v$ and $|R|_v=1$.  Globally
$\ha{R}\le\ha{a}+\ha{b}\le Ch$.  On the common special
fiber, relative Frobenius $F:\overline A\to
\overline A^{(p)}$ satisfies $F\alpha=\alpha^{(p)}F$
for every endomorphism $\alpha$, and similarly for the
centered endomorphisms.  Thus conjugation by $F$, which is
the normalized isometry on the operator system, carries
each labeled vector to its primed vector.  This records an
identity of actual specialized endomorphisms, in addition
to the de Rham equation.  No restriction on the Newton
polygon has been used.

For the local sums retain the selected embeddings of $M$ with their
absolute weights.  If a later algebraic operation introduces another
coefficient field $F'$, work over $MF'$ and take all extensions of those
selected embeddings.  The common equation, the equality $|R|_v=1$, and
the local norm bounds are preserved by restriction of absolute values,
and the new weights sum to the old ones.  This step requires no new
choice of a Frobenius lift.  Thus the field indexing the selected
embeddings can be large, while the field defining the algebraic tuple
has polynomial degree.  The degree of the normal closure never enters
the coordinate-degree, height, or grouping estimates.

There are in total polynomially many groups.  Apply
\Cref{lem:proximity} and select one to obtain
\eqref{eq:group-smallness}.  Put the ordinary coordinates
$z,P$ and, where present, $z',P',R$ in $U$.  The constants
$B,B'$ belong to fixed data; $T$ is a coefficient of the
equation, not an extra variable.  Equations
\eqref{eq:chart-relations} require the base coordinates to
be those of the selected inverse branches.  Their inclusion
keeps the number of variables and equations fixed and
prevents another sheet of the rational parameter from
being treated as the chosen germ.  Equation
\eqref{eq:common-size} follows from
\Cref{prop:dr-size}, its conjugate version, and the local
and global bounds for $R$.

\emph{Coefficient heights.}
It remains to bound the common series, uniformly across its selected
places.  Outside the fixed exceptional set, the algebraic chart series
and the connection matrix in that chart
have integral Taylor coefficients.  If
$Z(X)=\sum Z_jX^j$ is an identity-normalized horizontal
matrix with integral ordinary connection matrix, its
recursion shows inductively that
\[
 Z_j\in (j!)^{-1}\mathrm M_r(\Ocal_v).
\]
For example $(j+1)Z_{j+1}$ is a sum of $Z_k$ times
integral connection coefficients for $k\le j$, and
$j!/k!$ is integral.  The same holds for inverse
horizontal matrices and the fixed rank-five system.
Hence $\norm{Z_j}_v\le\exp(j\alpha_p)$ at a good
prime.  Integral algebraic chart series satisfy the
stronger bound $\norm{z_j}_v\le1$.

At every fixed exceptional finite place the implicit
function theorem gives a positive radius of convergence
for each algebraic chart series; its coefficients and
those of the connection are bounded by $c_v^{j+1}$ for
a fixed $c_v\ge1$.  The same recursion then bounds a
horizontal coefficient by $c_v^j|1/j!|_v$, after increasing
$c_v$.  At infinity the ordinary analytic solution and
the chart series have fixed positive convergence radii,
so their coefficients grow at most exponentially.
Since only finitely many $c_v$ differ from one, summing
these bounds and using
$\sum_{p}v_p(j!)\log p=\log(j!)$ gives, for example,
\[
 \ha{Z_0,\ldots,Z_j,z_0,\ldots,z_j}
       \le C(j+1)\log(j+2).
\]
Only this polynomial estimate for logarithmic heights is
needed.  In the second series $b=aR$, the coefficient of
$a^j$ has an additional factor $R^j$, explaining the
linear degree bound.  The remaining coefficients are
fixed algebraic coefficients, together with the four
rational constants $T$ in the constant case.  Their
affine heights, including scale, have just been bounded.
Polynomial operations of fixed degree and concatenation through order
$l$ preserve a polynomial height bound.  This proves
\eqref{eq:coefficient-size} for the complete coefficient tuple.

\emph{Local tails.}
At an unexceptional finite place put
$M_v=\max(1,\norm{U}_v)$.  The ultrametric inequality
and $|R|_v=1$ give, for a fixed integer $b$ and every $l$,
\[
 \left\|\sum_{j\ge l}a^j e_j(U)\right\|_v
       \le M_v^b |a|_v^l\exp(l\alpha_p).
\]
For $l=0$ the tail is the entire common relation and
vanishes at the target.  For $l\ge1$ the constant matrix
$T$ does not occur in the tail.
The displayed inequality follows by taking the maximum
over $j\ge l$; \eqref{eq:factorial-disk} makes those
bounds decrease.  There is no multiplicative constant
greater than one at every prime.  Fixed products or
tensor expressions satisfy the same estimate after
increasing $b$ and the fixed coefficient bounds, because
the total Taylor index controls the factorial bound.

At the finitely many exceptional places and at infinity,
choose the disks strictly inside a fixed convergence
disk.  Cauchy estimates, or the local differential
recursion, then give the same tail estimate with
$\alpha_p$ replaced by a fixed constant; at infinity
the geometric-series sum contributes one fixed further
constant.  Consequently one may take
\[
 b_v\le C\log M_v+C\alpha_p
\]
at the unexceptional primes, with fixed additional
constants at the other places.  The contribution of
$\alpha_p$ over all selected primes is
$O(1+\log^2(dh))$ by \eqref{eq:prime-support} and the
calculation in \Cref{lem:proximity}.  Summing and using
\eqref{eq:common-size} proves
\eqref{eq:arithmetic-tail}--\eqref{eq:tail-sum}.
All operations used involve a fixed number of fixed
ordinary systems and their fixed tensor constructions,
as asserted.
\end{proof}

The equations in \Cref{prop:common-relations} provide common
arithmetic relations, including at primes where the reduction is
supersingular.  Their nonidentity on the algebraic varieties
encountered by interpolation is a separate requirement, addressed
in the next section.
